% 第 6 章：三角函数
% 对应湘教版必修第一册 第 5 章
% 说明：本章按内容编排，课时划分由教师把握。
\section{三角函数}

\begin{learnobjective}
\begin{itemize}
  \item 理解任意角的概念，会进行角度与弧度的互化，会计算弧长与扇形面积；
  \item 会用单位圆定义任意角的正弦、余弦、正切，会求特殊角的三角函数值；
  \item 掌握同角三角函数关系与诱导公式，会化简与求值；
  \item 会画正弦、余弦、正切的图象，说出它们的周期性、奇偶性与单调性；
  \item 理解 $y=A\sin(\omega x+\varphi)$ 中 $A,\omega,\varphi$ 的意义，会作图象变换；
  \item 会用三角函数建立简单的周期现象模型。
\end{itemize}
\end{learnobjective}

\begin{preview}
\subsection*{知识回顾}
\begin{itemize}
  \item 锐角三角函数：在直角三角形中，$\sin\alpha=$ \blank，$\cos\alpha=$ \blank；
  \item $30^\circ$、$45^\circ$、$60^\circ$ 的正弦值分别是 \blank、\blank、\blank；
  \item 半径为 $r$ 的圆，周长是 \blank。
\end{itemize}

\subsection*{教材速读}
阅读本章内容，完成下列问题：
\begin{enumerate}
  \item $180^\circ=$ \blank rad，$1^\circ=$ \blank rad；
  \item 弧度制下弧长公式是 $l=$ \blank；
  \item $y=\sin x$ 的最小正周期是 \blank，$y=\tan x$ 的最小正周期是 \blank。
\end{enumerate}
\end{preview}

\begin{explore}

\subsection{任意角与弧度制}

\key{探究一：任意角与终边相同的角}
\begin{itemize}
  \item 角可以看成平面内一条射线绕着端点旋转所形成的图形。
        \key{逆时针}旋转成\key{正角}，\key{顺时针}旋转成\key{负角}，没有旋转成\key{零角}。
  \item 这样，角的范围就不再是 $0^\circ\sim180^\circ$，而是任意大小（可以超过 $360^\circ$，也可以是负的）。
  \item \key{终边相同的角}：所有与角 $\alpha$ 终边相同的角，连同 $\alpha$ 在内，可写成
        \[
          \{\beta\mid\beta=\alpha+k\cdot360^\circ,\ k\in\mathbb{Z}\}.
        \]
        即相差 $360^\circ$ 的整数倍。
  \item 判断象限：把角化到 $0^\circ\sim360^\circ$ 之间再看终边落在哪个象限。
\end{itemize}

\begin{example}
(1) 写出与 $30^\circ$ 终边相同的一个负角；
(2) 在 $0^\circ\sim360^\circ$ 内，找出与 $-400^\circ$ 终边相同的角；
(3) 判断 $-120^\circ$ 是第几象限角。
\end{example}
\begin{solution}
(1) $30^\circ-360^\circ=-330^\circ$（答案不唯一）；
(2) $-400^\circ=-360^\circ-40^\circ$，终边与 $-40^\circ$ 相同，再加 $360^\circ$ 得 $320^\circ$；
(3) $-120^\circ$ 的终边与 $240^\circ$ 相同，在第三象限。
\end{solution}

\key{探究二：弧度制}
\begin{itemize}
  \item 把长度等于\key{半径}长的弧所对的圆心角叫做 $1$ 弧度的角，记作 $1\ \text{rad}$。
  \item 换算：$180^\circ=\pi\ \text{rad}$，$1^\circ=\dfrac{\pi}{180}\ \text{rad}$，
        $1\ \text{rad}=\left(\dfrac{180}{\pi}\right)^\circ\approx57.3^\circ$。
  \item 弧长公式：$l=|\alpha|\,r$；扇形面积：$S=\dfrac12lr=\dfrac12|\alpha|\,r^2$。
        （$\alpha$ 用弧度）
\end{itemize}

\begin{example}
把 $60^\circ$、$-135^\circ$ 化成弧度，把 $\dfrac{2\pi}{3}$ 化成角度。
\end{example}
\begin{solution}
$60^\circ=60\times\dfrac{\pi}{180}=\dfrac{\pi}{3}$；
$-135^\circ=-135\times\dfrac{\pi}{180}=-\dfrac{3\pi}{4}$；
$\dfrac{2\pi}{3}=\dfrac{2\pi}{3}\times\dfrac{180}{\pi}=120^\circ$。
\end{solution}

\begin{problem}
已知扇形的半径为 $2$，圆心角为 $\dfrac{\pi}{3}$，求弧长与面积。
\end{problem}
\begin{solution}
$l=|\alpha|r=\dfrac{\pi}{3}\times2=\dfrac{2\pi}{3}$；
$S=\dfrac12lr=\dfrac12\times\dfrac{2\pi}{3}\times2=\dfrac{2\pi}{3}$。
\end{solution}

\key{为什么用弧度}
角度制把圆周分成 $360$ 份，是人为约定；弧度用“弧长 $\div$ 半径”定义，
与半径无关，公式更简洁（$l=\alpha r$、$S=\frac12\alpha r^2$），
也使 $\sin x$、$\cos x$ 的导数公式变得干净。因此\key{高等数学与物理中一律用弧度}。

\subsection{任意角的三角函数}

\key{探究三：三角函数的定义}
\begin{itemize}
  \item 取\key{单位圆}（圆心在原点、半径为 $1$）。设角 $\alpha$ 的终边与单位圆交于 $P(x,y)$，则
        \[
          \sin\alpha=y,\qquad \cos\alpha=x,\qquad \tan\alpha=\frac{y}{x}\ (x\neq0).
        \]
  \item 更一般地，若角 $\alpha$ 终边过点 $P(x,y)$，记 $r=\sqrt{x^2+y^2}$，则
        $\sin\alpha=\dfrac{y}{r}$，$\cos\alpha=\dfrac{x}{r}$，$\tan\alpha=\dfrac{y}{x}$。
  \item \key{三角函数线}：在单位圆中，有向线段 $MP$ 是正弦线，$OM$ 是余弦线，
        过点 $(1,0)$ 作切线，$AT$ 是正切线。
\end{itemize}

\begin{example}
已知角 $\alpha$ 的终边过点 $(3,4)$，求 $\sin\alpha$、$\cos\alpha$、$\tan\alpha$；
再求角 $\dfrac{\pi}{3}$ 的三个三角函数值。
\end{example}
\begin{solution}
$r=\sqrt{3^2+4^2}=5$，故 $\sin\alpha=\dfrac45$，$\cos\alpha=\dfrac35$，$\tan\alpha=\dfrac43$。
$\dfrac{\pi}{3}$ 的终边与单位圆交于 $\left(\dfrac12,\dfrac{\sqrt3}{2}\right)$，
故 $\sin\dfrac{\pi}{3}=\dfrac{\sqrt3}{2}$，$\cos\dfrac{\pi}{3}=\dfrac12$，$\tan\dfrac{\pi}{3}=\sqrt3$。
\end{solution}

\key{探究四：同角关系与诱导公式}
\begin{itemize}
  \item \key{同角关系}：$\sin^2\alpha+\cos^2\alpha=1$，$\tan\alpha=\dfrac{\sin\alpha}{\cos\alpha}$。
  \item \key{诱导公式}（把任意角化成锐角）：
        \[
          \begin{aligned}
            \sin(-\alpha)&=-\sin\alpha, & \cos(-\alpha)&=\cos\alpha,\\
            \sin(\pi-\alpha)&=\sin\alpha, & \cos(\pi-\alpha)&=-\cos\alpha,\\
            \sin(\pi+\alpha)&=-\sin\alpha, & \cos(\pi+\alpha)&=-\cos\alpha,\\
            \sin\!\left(\frac{\pi}{2}-\alpha\right)&=\cos\alpha, &
            \cos\!\left(\frac{\pi}{2}-\alpha\right)&=\sin\alpha.
          \end{aligned}
        \]
  \item 口诀：\key{奇变偶不变，符号看象限}（把 $\alpha$ 当锐角，看变换后的角落在哪个象限）。
\end{itemize}

\begin{example}
已知 $\sin\alpha=\dfrac35$，且 $\alpha$ 是第二象限角，求 $\cos\alpha$ 与 $\tan\alpha$。
\end{example}
\begin{solution}
由 $\sin^2\alpha+\cos^2\alpha=1$ 得 $\cos^2\alpha=1-\dfrac{9}{25}=\dfrac{16}{25}$，
故 $\cos\alpha=\pm\dfrac45$。$\alpha$ 在第二象限，$\cos\alpha<0$，取 $\cos\alpha=-\dfrac45$。
$\tan\alpha=\dfrac{\sin\alpha}{\cos\alpha}=-\dfrac34$。
\end{solution}

\begin{problem}
化简：$\sin(180^\circ+\alpha)\cos(-\alpha)$；并求 $\sin\dfrac{7\pi}{6}$。
\end{problem}
\begin{solution}
$\sin(180^\circ+\alpha)=-\sin\alpha$，$\cos(-\alpha)=\cos\alpha$，故原式 $=-\sin\alpha\cos\alpha$。
$\sin\dfrac{7\pi}{6}=\sin\left(\pi+\dfrac{\pi}{6}\right)=-\sin\dfrac{\pi}{6}=-\dfrac12$。
\end{solution}

\subsection{三角函数的图象与性质}

\key{探究五：正弦与余弦的图象和性质}
\begin{itemize}
  \item 用单位圆描点，或用\key{五点法}：取 $x=0,\ \dfrac{\pi}{2},\ \pi,\ \dfrac{3\pi}{2},\ 2\pi$
        五个关键点连线。
  \item $y=\sin x$：定义域 $\mathbb{R}$，值域 $[-1,1]$，最小正周期 $2\pi$，奇函数；
        在 $\left[-\dfrac{\pi}{2}+2k\pi,\dfrac{\pi}{2}+2k\pi\right]$ 上递增，
        在 $\left[\dfrac{\pi}{2}+2k\pi,\dfrac{3\pi}{2}+2k\pi\right]$ 上递减。
  \item $y=\cos x$：定义域 $\mathbb{R}$，值域 $[-1,1]$，最小正周期 $2\pi$，偶函数；
        在 $[2k\pi,\pi+2k\pi]$ 上递减，在 $[\pi+2k\pi,2\pi+2k\pi]$ 上递增。
\end{itemize}

\begin{figure}[htbp]
  \centering
  \begin{tikzpicture}
    \begin{axis}[
      width=0.80\textwidth, height=6.4cm,
      xmin=-2*pi, xmax=2*pi, ymin=-1.6, ymax=1.6,
      xtick={-2*pi, -3*pi/2, -pi, -pi/2, 0, pi/2, pi, 3*pi/2, 2*pi},
      xticklabels={$-2\pi$,$-\frac{3\pi}{2}$,$-\pi$,$-\frac{\pi}{2}$,$0$,$\frac{\pi}{2}$,$\pi$,$\frac{3\pi}{2}$,$2\pi$},
      trig format plots=rad,
      xlabel=$x$, ylabel=$y$,
      grid=major, grid style={dashed,gray!35},
      legend pos=north east, legend style={font=\footnotesize},
      tick label style={font=\footnotesize},
      title={\small $y=\sin x$ 与 $y=\cos x$（周期 $2\pi$）},
      title style={font=\small},
    ]
      \addplot[domain=-2*pi:2*pi, samples=240, thick, main!80!black] {sin(x)};
        \addlegendentry{$y=\sin x$}
      \addplot[domain=-2*pi:2*pi, samples=240, thick, second] {cos(x)};
        \addlegendentry{$y=\cos x$}
      \addplot[only marks, mark=*, mark size=1.4pt, black]
        coordinates {(0,0) (pi/2,1) (pi,0) (3*pi/2,-1) (2*pi,0)};
    \end{axis}
  \end{tikzpicture}
  \caption{正弦（过原点的实线）与余弦：黑点是 $y=\sin x$ 的五点法关键点}
\end{figure}

\key{探究六：正切的图象和性质}
\begin{itemize}
  \item $y=\tan x$：定义域 $\left\{x\mid x\neq\dfrac{\pi}{2}+k\pi,\ k\in\mathbb{Z}\right\}$，
        值域 $\mathbb{R}$，最小正周期 $\pi$，奇函数。
  \item 它在每个区间 $\left(-\dfrac{\pi}{2}+k\pi,\dfrac{\pi}{2}+k\pi\right)$ 内单调递增。
  \item 直线 $x=\dfrac{\pi}{2}+k\pi$ 是图象的\key{渐近线}：函数值趋向无穷，但永远取不到。
\end{itemize}

\begin{figure}[htbp]
  \centering
  \begin{tikzpicture}
    \begin{axis}[
      width=0.80\textwidth, height=6.4cm,
      xmin=-4.8, xmax=4.8, ymin=-4.5, ymax=4.5,
      xtick={-3*pi/2, -pi, -pi/2, 0, pi/2, pi, 3*pi/2},
      xticklabels={$-\frac{3\pi}{2}$,$-\pi$,$-\frac{\pi}{2}$,$0$,$\frac{\pi}{2}$,$\pi$,$\frac{3\pi}{2}$},
      trig format plots=rad,
      xlabel=$x$, ylabel=$y$,
      grid=major, grid style={dashed,gray!35},
      legend pos=north west, legend style={font=\footnotesize},
      tick label style={font=\footnotesize},
      title={\small $y=\tan x$：周期 $\pi$，虚线为渐近线},
      title style={font=\small},
    ]
      \addplot[domain=-1.32:1.32, samples=120, thick, main!80!black] {tan(x)};
      \addplot[domain=1.83:4.45, samples=120, thick, main!80!black] {tan(x)};
      \addplot[domain=-4.45:-1.83, samples=120, thick, main!80!black] {tan(x)};
      \addplot[dashed, gray] coordinates {(-pi/2,-4.5) (-pi/2,4.5)};
      \addplot[dashed, gray] coordinates {(pi/2,-4.5) (pi/2,4.5)};
      \addplot[dashed, gray] coordinates {(-3*pi/2,-4.5) (-3*pi/2,4.5)};
      \addplot[dashed, gray] coordinates {(3*pi/2,-4.5) (3*pi/2,4.5)};
      \addlegendentry{$y=\tan x$}
    \end{axis}
  \end{tikzpicture}
  \caption{正切曲线：每个周期内递增，渐近线处断开}
\end{figure}

\subsection{函数 $y=A\sin(\omega x+\varphi)$ 的图象与性质}

\key{探究七：三个参数的作用}
\begin{itemize}
  \item 形如 $y=A\sin(\omega x+\varphi)$（$A>0$，$\omega>0$）的函数：
        \begin{itemize}
          \item $|A|$ 叫\key{振幅}：图象纵向拉伸（$|A|>1$）或压缩（$|A|<1$）；
          \item $\omega$ 决定\key{周期} $T=\dfrac{2\pi}{\omega}$：横向压缩或拉伸；
          \item $\varphi$ 叫\key{初相}：决定左右平移，图象左移 $\dfrac{\varphi}{\omega}$ 个单位。
        \end{itemize}
  \item 图象的画法：先画 $y=\sin x$，再按“平移 $\to$ 横缩 $\to$ 纵伸”三步变换；
        或代入五个关键点（令 $\omega x+\varphi$ 依次等于 $0,\frac{\pi}{2},\pi,\frac{3\pi}{2},2\pi$）。
\end{itemize}

\begin{example}
求 $y=2\sin\left(2x+\dfrac{\pi}{3}\right)$ 的振幅、周期与初相，并说出它可由 $y=\sin x$ 怎样变换得到。
\end{example}
\begin{solution}
振幅 $A=2$，周期 $T=\dfrac{2\pi}{2}=\pi$，初相 $\varphi=\dfrac{\pi}{3}$。
变换：$y=\sin x$ 向左平移 $\dfrac{\pi}{3}$ 个单位得 $y=\sin\left(x+\dfrac{\pi}{3}\right)$；
横坐标缩短为原来的 $\dfrac12$ 得 $y=\sin\left(2x+\dfrac{\pi}{3}\right)$；
纵坐标伸长为原来的 $2$ 倍得 $y=2\sin\left(2x+\dfrac{\pi}{3}\right)$。
\end{solution}

\begin{figure}[htbp]
  \centering
  \begin{tikzpicture}
    \begin{axis}[
      width=0.80\textwidth, height=6.4cm,
      xmin=0, xmax=2*pi, ymin=-2.6, ymax=2.6,
      xtick={0, pi/2, pi, 3*pi/2, 2*pi},
      xticklabels={$0$,$\frac{\pi}{2}$,$\pi$,$\frac{3\pi}{2}$,$2\pi$},
      trig format plots=rad,
      xlabel=$x$, ylabel=$y$,
      grid=major, grid style={dashed,gray!35},
      legend pos=north east, legend style={font=\footnotesize},
      tick label style={font=\footnotesize},
      title={\small $y=\sin x$ 与 $y=2\sin\left(2x+\frac{\pi}{3}\right)$},
      title style={font=\small},
    ]
      \addplot[domain=0:2*pi, samples=200, thick, gray] {sin(x)};
        \addlegendentry{$y=\sin x$}
      \addplot[domain=0:2*pi, samples=400, thick, main!80!black] {2*sin(2*x+pi/3)};
        \addlegendentry{$y=2\sin\left(2x+\frac{\pi}{3}\right)$}
    \end{axis}
  \end{tikzpicture}
  \caption{振幅变大、周期变小、整体左移：同一族曲线形状相同，位置与大小不同}
\end{figure}

\begin{problem}
函数 $y=3\sin\left(\dfrac{x}{2}-\dfrac{\pi}{4}\right)$ 的振幅、周期、初相各是多少？
它的图象可由 $y=\sin x$ 怎样变换得到？
\end{problem}
\begin{solution}
振幅 $3$，周期 $T=\dfrac{2\pi}{1/2}=4\pi$，初相 $-\dfrac{\pi}{4}$。
先把 $y=\sin x$ 右移 $\dfrac{\pi}{4}$ 得 $y=\sin\left(x-\dfrac{\pi}{4}\right)$，
再把横坐标伸长为原来的 $2$ 倍得 $y=\sin\left(\dfrac{x}{2}-\dfrac{\pi}{4}\right)$，
最后把纵坐标伸长为原来的 $3$ 倍。
\end{solution}

\subsection{三角函数模型的简单应用}

\key{探究八：周期现象的刻画}
\begin{itemize}
  \item 很多现象每过相同时间就重复一次：昼夜、四季、心跳、钟摆、交流电。
        这类现象可用 $y=A\sin(\omega x+\varphi)+b$ 刻画。
  \item 其中 $b$ 是\key{平衡位置}（平均值），$A$ 是偏离平衡位置的最大幅度。
\end{itemize}

\begin{example}
某地一天的温度可用 $T(t)=8\sin\left(\dfrac{\pi}{12}t-\dfrac{7\pi}{12}\right)+20$（$t$ 为小时）近似描述。
求这一天的最高温度、最低温度与周期。
\end{example}
\begin{solution}
振幅 $A=8$，平衡位置 $b=20$，故最高 $20+8=28$（℃），最低 $20-8=12$（℃）。
周期 $T=\dfrac{2\pi}{\pi/12}=24$（小时），正好是一天。
\end{solution}

\begin{problem}
一个弹簧振子做简谐运动，位移 $y$（cm）与时间 $t$（s）满足 $y=5\sin(2\pi t)$。
求振幅、周期，并求 $t=0.25$ s 时的位移。
\end{problem}
\begin{solution}
振幅 $5$ cm，周期 $T=\dfrac{2\pi}{2\pi}=1$ s。
$t=0.25$ 时 $y=5\sin\dfrac{\pi}{2}=5$ cm，即到达正向最大位移处。
\end{solution}

\key{数学文化：三角学的历史}
三角学起源于测量。古希腊的喜帕恰斯、托勒密集古埃及与巴比伦的成果，
编出最早的弦表，用于天文计算；印度人引入“正弦”概念并给出正弦表；
阿拉伯学者把三角学系统化，并传入欧洲。17 世纪以后，三角函数成为
分析学的一部分（傅里叶级数说明：任何周期现象都可以用三角函数叠加来逼近）。
今天三角函数不仅是解题工具，更是描述\key{周期现象}的通用语言。

\end{explore}

\begin{inclasstest}
\begin{enumerate}
  \item 把 $150^\circ$、$-210^\circ$ 化成弧度；把 $\dfrac{5\pi}{6}$ 化成角度；
  \item 写出与 $-60^\circ$ 终边相同的角的集合，并找出在 $-360^\circ\sim0^\circ$ 内的那个角；
  \item 已知角 $\alpha$ 的终边过点 $(-4,3)$，求 $\sin\alpha$、$\cos\alpha$、$\tan\alpha$；
  \item 已知 $\cos\alpha=-\dfrac{5}{13}$，且 $\alpha$ 在第三象限，求 $\sin\alpha$、$\tan\alpha$；
  \item 化简 $\cos\left(\dfrac{\pi}{2}+\alpha\right)+\sin(\pi-\alpha)$；
  \item 求 $y=\sin x$ 与 $y=\cos x$ 的单调递增区间；
  \item 求 $y=\tan x$ 的定义域与周期；
  \item 求 $y=4\sin\left(3x-\dfrac{\pi}{6}\right)$ 的振幅、周期与初相。
\end{enumerate}
\end{inclasstest}

\begin{afterclass}
\begin{enumerate}
  \item 把 $22^\circ30'$ 化成弧度；把 $\dfrac{3\pi}{4}$ rad 化成角度；
  \item 半径为 $6$、圆心角为 $120^\circ$ 的扇形，求弧长与面积；
  \item 已知 $\sin\alpha=\dfrac{12}{13}$，$\alpha$ 为第一象限角，求 $\cos\alpha$、$\tan\alpha$；
  \item 化简 $\dfrac{\sin(2\pi-\alpha)\cos(\pi+\alpha)}{\tan(-\alpha)}$；
  \item 用五点法在区间 $[0,2\pi]$ 上画出 $y=1+2\sin x$ 的图象（写出五个关键点）；
  \item 判断 $y=\sin x+\cos x$ 的奇偶性，并求它的最大值；
  \item 函数 $y=A\sin(\omega x+\varphi)$ 的图象相邻两条最高点距离为 $\pi$，且过点 $\left(0,\dfrac12\right)$、振幅为 $1$，写出它的一个解析式；
  \item 阅读“数学文化：三角学的历史”，举例说明三角函数在生活中的一个周期现象。
\end{enumerate}
\reading{拓展阅读 06}{三角与周期}
\end{afterclass}
